%% kapa: exam=1 round=23 year=2012 period=1 form=A subject=civil_law subject_ko=민법 seq=22 number=22
%% source: Q-Net 감정평가사 기출문제 (cst00309) artlSeq=5052578 file=raw/1008/5052578/2012년 제23회 감정평가사 제1차시험 기출문제(1교시 A형).pdf
%% conversion: kordoc 4.24.51 -> tools/gen1.py (auto); review=proofread-sonnet; flags=none
\begin{question}{22}
법률행위의 해석에 관한 설명으로 옳지 않은 것은? (다툼이 있으면 판례에 의함)
\begin{choices}
  \item 계약의 해석은 당사자가 그 표시행위에 부여한 객관적 의미를 명백하게 확정하는 것이다.
  \item 계약의 당사자가 누구인지는 그 계약에 관여한 당사자의 의사해석의 문제이다.
  \item 진의 아닌 의사표시에 있어서 진의는 표의자가 진정으로 마음 속에서 바라는 사항을 의미한다.
  \item 계약의 해석을 통하여 보충되는 당사자의 의사는 당사자의 주관적 의사가 아니라 거래관행이나 신의칙 등에 의하여 객관적으로 추인되는 의사를 의미한다.
  \item 타인 명의로 계약을 체결하는 경우 계약을 체결한 자와 상대방의 의사가 일치한다면 그 일치된 의사대로 당사자를 확정하면 된다.
\end{choices}
\end{question}
\answer{3}
